\chapter{Usability Testing Documentation}
\label{app:usability-testing}

This appendix presents the questions, task instructions and observation procedure used in the two Omnisound usability tests. Selected individual records show how participants approached the tasks and what they said about the interface. The full results and their interpretation are presented in Chapters 5 and 6.

\section{Test Preparation}

The group tested two interactive versions of Omnisound. Three students, P1--P3, used the first prototype, V1. Three different students, P4--P6, used the revised prototype, V2. Both rounds used a laptop with the prototype running in a web browser.

Yacoub Dawli moderated the sessions and presented the tasks. Obid Alrhman Aljabi recorded participant behaviour and comments. Sarshar Chavoshi Rad recorded completion times, clicks and navigation routes.

The tests examined effectiveness through task completion, efficiency through time and navigation, and learnability through participants' understanding of the interface on first use. The group also looked at the visibility of important destinations and the feedback given after actions such as saving an item.

Before testing, Yacoub introduced the group as students at Mittuniversitetet working on Omnisound. He explained that the interface was being tested, not the participant. Participation was voluntary, and participants could stop at any time or raise questions and concerns. The group would observe their use of the prototype and take anonymous written notes for the university project; names would not appear in the report. Participants were encouraged to think aloud about what they were looking for, expected to happen or found unclear. The moderator then asked whether they were comfortable participating. The records include agreement before the tasks began.

\section{Pre-Test Questions}

The background questions covered previous experience with audio applications, listening habits and usual devices. These helped the group understand participants' familiarity with the kinds of content and activities supported by Omnisound.

The following questions use the recorded V2 wording. V1 covered the same topics with slightly different wording.

\begin{enumerate}
  \setlength{\itemsep}{2pt}
  \setlength{\parsep}{0pt}
  \item Which audio applications do you normally use?
  \item Do you mainly listen to music, podcasts, audiobooks, or a mixture?
  \item How often do you use audio applications?
  \item Do you mainly use them on mobile, desktop, or both?
\end{enumerate}

\section{Task Instructions}

The moderator presented one task at a time. The instructions below are the main V1 prompts. They described the goal without telling participants which control or route to use.

\begin{samepage}
\subsection*{Task 1: Find and Save Content}

\begin{quote}
``You want to listen to Weezer later. Find Weezer and save it somewhere so you can easily find it again.''
\end{quote}

\noindent\textbf{Success:} Weezer is saved and can be located again in Library.
\end{samepage}

\begin{samepage}
\subsection*{Task 2: Resume Unfinished Audio}

\begin{quote}
``Earlier you started listening to an audiobook but did not finish it. Find your unfinished audiobook and continue from where you stopped.''
\end{quote}

\noindent\textbf{Success:} The participant identifies the unfinished audiobook, selects Resume and reaches the player at the existing listening position.
\end{samepage}

\begin{samepage}
\subsection*{Task 3: Browse Podcasts}

\begin{quote}
``You are only interested in podcasts right now. Find where you can browse podcasts.''
\end{quote}

\noindent\textbf{Success:} The participant reaches Podcasts.
\end{samepage}

V2 used slightly different wording while preserving these three task goals and completion criteria. The unfinished-audiobook prompt did not name Continue Listening, leaving participants free to choose their own route.

\section{Think-Aloud Procedure}

Participants were encouraged to explain their thoughts while working through the tasks. When prompting was needed, the moderator used neutral questions:

\begin{itemize}
  \setlength{\itemsep}{2pt}
  \setlength{\parsep}{0pt}
  \item ``What are you thinking?''
  \item ``What are you looking for?''
  \item ``What did you expect to happen?''
  \item ``What do you think this control does?''
\end{itemize}

When participants asked what to select, they were encouraged to try what made sense to them rather than being immediately directed to a correct button. The recorded response was: ``Try what you think makes sense.'' This allowed the group to observe their choices and any difficulties.

For each task, the observers recorded task success, time in seconds, click count, navigation route, hesitation, incorrect actions, assistance and participant comments. These notes helped explain how the task was completed, rather than relying only on its completion time.

\section{Post-Test Questions}

After the tasks, participants were asked about their overall experience and any changes they would suggest. The wording differed slightly between participants. The following full-sentence questions represent the main V1 topics.

\begin{enumerate}
  \setlength{\itemsep}{2pt}
  \setlength{\parsep}{0pt}
  \item What was your overall impression?
  \item What felt easiest?
  \item What was most confusing?
  \item Was there anything you expected but could not find?
  \item Were any labels confusing?
  \item What would you change first?
  \item Was anything important missing?
\end{enumerate}

The V2 questions covered overall impressions, the easiest and most confusing parts, unclear controls and labels, discoverability and suggested changes. Open questions also brought out comments about navigation and functions participants wanted. Recorded V2 questions included:

\begin{enumerate}
  \setlength{\itemsep}{2pt}
  \setlength{\parsep}{0pt}
  \item ``Overall impression?''
  \item ``Easiest part?''
  \item ``Most confusing?''
  \item ``Anything unclear?''
  \item ``What would you change?''
  \item ``Did anything feel difficult to discover?''
  \item ``Was Save clear?''
\end{enumerate}

The last two questions addressed discoverability and the clarity of Save directly. Participants could discuss Continue Listening in their answers without the moderator naming it in those questions. Suggestions for missing functions included Recently Played and Library sorting.

\section{Selected V1 Observations}

P1--P3 completed all three tasks without direct help. Table~\ref{tab:appendix-v1} records their task times, click counts and audiobook routes. Each task cell gives the time in seconds followed by the number of clicks.

\begin{table}[!htbp]
  \centering
  \caption{V1 individual task times and click counts.}
  \label{tab:appendix-v1}
  \begingroup
  \small
  \renewcommand{\arraystretch}{1.3}
  \begin{tabular}{@{}lcccp{0.24\textwidth}@{}}
    \hline
    \textbf{Participant} & \textbf{Task 1} & \textbf{Task 2} & \textbf{Task 3} & \textbf{Audiobook route} \\
    \hline
    P1 & 31 s / 5 clicks & 19 s / 2 clicks & 4 s / 1 click & Continue Listening \\
    P2 & 24 s / 5 clicks & 27 s / 3 clicks & 5 s / 1 click & Audiobooks \\
    P3 & 28 s / 4 clicks & 35 s / 4 clicks & 7 s / 2 clicks & Library \\
    \hline
  \end{tabular}
  \endgroup
\end{table}

The following post-test comments explain the main concerns:

\begin{itemize}
  \setlength{\itemsep}{2pt}
  \setlength{\parsep}{0pt}
  \item \textbf{P1:} ``Maybe make Save slightly more obvious.''
  \item \textbf{P2:} ``I didn't immediately notice Continue Listening. I went to Audiobooks instead.''
  \item \textbf{P3:} ``Maybe make Continue Listening more obvious. The big banner above it takes a lot of space.''
\end{itemize}

P1 briefly hesitated when locating Save beside Play. P2 and P3 completed the audiobook task through other destinations after overlooking Continue Listening. P3 also selected Music before correcting the choice to Podcasts without help. These observations informed the priorities for V2: a smaller Home banner, a more noticeable Continue Listening section and a clearer Save/Saved control.

\section{Selected V2 Observations}

P4--P6 completed all three tasks without direct help. Table~\ref{tab:appendix-v2} uses the same format as the V1 records.

\begin{table}[!htbp]
  \centering
  \caption{V2 individual task times and click counts.}
  \label{tab:appendix-v2}
  \begingroup
  \small
  \renewcommand{\arraystretch}{1.3}
  \begin{tabular}{@{}lcccp{0.24\textwidth}@{}}
    \hline
    \textbf{Participant} & \textbf{Task 1} & \textbf{Task 2} & \textbf{Task 3} & \textbf{Audiobook route} \\
    \hline
    P4 & 22 s / 5 clicks & 13 s / 2 clicks & 4 s / 1 click & Continue Listening \\
    P5 & 24 s / 5 clicks & 15 s / 2 clicks & 4 s / 1 click & Continue Listening \\
    P6 & 21 s / 4 clicks & 12 s / 2 clicks & 3 s / 1 click & Continue Listening \\
    \hline
  \end{tabular}
  \endgroup
\end{table}

Selected comments concerned saving feedback and a possible later feature:

\begin{itemize}
  \setlength{\itemsep}{2pt}
  \setlength{\parsep}{0pt}
  \item \textbf{P4:} ``Yes. The heart and word Save made it obvious.''
  \item \textbf{P5:} ``Okay, now I know it worked.''
  \item \textbf{P6:} ``Maybe let me sort the Library, but that's not really necessary here.''
\end{itemize}

P5 paused for approximately one second while looking between Play and Save, then saved the item and recognised the Saved state. All three participants used Continue Listening for the audiobook task, with two clicks and no recorded hesitation. No incorrect category selections were recorded for the podcast task.

The two rounds involved only three participants each and used different groups. These records describe the tested interactions and do not establish that the redesign alone caused differences between the rounds.
